\chapter{Quantum Mechanics without Ontology? The Arena Problem}\label{ch:ontology}
\index{ontology}\index{epistemology}\index{configuration space}\index{arena problem}\index{PBR theorem}\index{wave-function realism}\index{primitive ontology}

\begin{quote}
\itshape Whatever the multi-dimensional wave function is claimed to be, RealQM is a different theory --- and that
is the point.\\
\normalfont\hfill --- working note, 2026.
\end{quote}
\bigskip

The companion Chapter~\ref{ch:philosophy} argued the general point --- that physics is, before it is anything
else, an account of what exists, and that a theory which will not say so pays for the omission in the measurement
problem and the century of interpretations. This chapter makes the diagnosis precise, and does so in a way
that sidesteps a debate it need not enter. Much energy has gone into deciding whether the wave function is
\emph{real} (ontic) or a summary of \emph{knowledge} (epistemic). For the purpose of an alternative like RealQM,
\emph{it does not matter}. Whatever the multi-dimensional $\Psi(\mathbf x_1,\dots,\mathbf x_N)$ is claimed to be,
RealQM is a \emph{fundamentally different} theory: its primary object is not $\Psi$ on configuration space but
charge in ordinary three-dimensional space. The metaphysics of the configuration-space wave function is a
question internal to the standard framework; RealQM steps outside that framework and is to be judged on its own
predictions and explanatory power. What can be said is that almost every foundational difficulty of the standard
theory flows from its choosing a fundamental object outside three-space --- a choice made, and contested, in the
spring of 1926, and a contingent turn rather than a forced one.

\section{Why the status of $\Psi$ does not decide the matter}\label{sec:onto-pbr}

Within the standard framework there is a genuine and technical debate about the status of the wave function.
Harrigan and Spekkens (2010) distinguished $\psi$-epistemic from $\psi$-ontic models, and the
Pusey--Barrett--Rudolph (PBR) theorem (2012) argued that, under a preparation-independence assumption,
$\psi$-epistemic models reproducing the quantum predictions cannot succeed --- so that, \emph{within that
framework}, the state is naturally taken to be ontic. This is an important result, but it is a result about
\emph{models of the standard formalism}: it assumes an underlying physical state and a model that reproduces
standard quantum mechanics, and it constrains how the standard $\Psi$ may be read. It does not, and cannot, compel
a \emph{different} theory to adopt $\Psi$ as its object.

RealQM is such a different theory. It does not offer a reading of the configuration-space formalism, ontic or
epistemic; it offers another fundamental object --- charge in three-dimensional space --- and treats the
multi-dimensional $\Psi$, if it uses it at all, as a derived, statistical description. Whether the standard
community settles on ontic or epistemic is therefore, for RealQM, beside the point: the question is not what
$\Psi$ \emph{is}, but whether a theory built on a different object, one already in three dimensions, can do the
physics. That RealQM stands aside from the $\Psi$-status dispute is not evasion but the natural consequence of not
sharing the standard theory's fundamental object; and Kohn's observation that this $\Psi$ is not even a recordable
thing beyond a modest number of particles only underlines how little is lost by declining to build one's ontology
on it.

\section{The arena problem}\label{sec:arena-problem}

Chapter~\ref{ch:philosophy} traced the measurement problem, the Heisenberg cut, and the nonlocality puzzles to a
single source, and that list need not be re-drawn. The compact statement is this: the standard theory's
fundamental object is a field on a space whose dimension grows with particle number and which does not inhabit the
three-dimensional space where instruments read definite facts, and the familiar pathologies are what any such
theory must incur. The link from the many-dimensional object to definite three-dimensional outcomes is not
supplied by the unitary dynamics and is put back by hand as collapse; whether one then calls the many-dimensional
object ontic or epistemic changes none of this. What RealQM contributes is not a verdict on that object but an
object of the other kind --- a fundamental description already in three-dimensional space, from which
three-dimensional facts need no recovery, because they are what the theory is about. That is the sense in which
the difficulty is one of \emph{arena}: not a real thing placed in the wrong space, but a whole theory built on an
object outside the space where its phenomena occur.

\section{Wave-function realism versus primitive ontology}\label{sec:wfr-po}

The contemporary metaphysics of quantum mechanics is, in effect, an argument over the arena. \emph{Wave-function
realism} (Albert; Ney; and, for the Everettian version, Wallace) accepts the configuration-space object as
fundamental and takes the high-dimensional space to be the real arena, our three-dimensional experience being
derivative; it buys realism at the price of demoting ordinary space to appearance, and owes an account of how the
manifest spatial world is recovered. \emph{Primitive-ontology} approaches (Allori; D\"urr--Goldstein--Zangh\`i;
Maudlin; and in spirit Bell) take the opposite route: they posit fundamental beables in three-dimensional space
--- particles, a matter-density field, or flash events --- and give the wave function a nomological or guiding
role rather than the role of the primary substance. On these views the wave function may be ontic in PBR's sense
while not being the \emph{primitive} ontology; what is fundamentally real is in three dimensions. The choice
between the two families is not settled by PBR, which speaks to the reality of the state but not to its arena; it
turns on which program better discharges the duties of a physical theory --- recovering the three-dimensional
phenomena and dissolving the measurement problem without ad hoc additions.

\section{The choice of 1926}\label{sec:onto-1926}

The arena was chosen quickly, and against protest. In the first months of 1926 Schr\"odinger's wave mechanics
represented the one-electron atom by a field $\psi(\mathbf x)$ on ordinary space, and he hoped to read
$|\psi|^2$ as a charge density there. The extension to $N$ interacting electrons required a single function
$\Psi(\mathbf x_1,\dots,\mathbf x_N)$ on configuration space; Born's statistical interpretation then read
$|\Psi|^2$ as a probability over configurations. Schr\"odinger, in his correspondence with Lorentz, and Einstein,
through the Solvay debates and later the EPR argument of 1935, both objected that a theory whose fundamental
object was a many-dimensional amplitude had given up a physical picture in three-space; both were overruled as the
G\"ottingen--Copenhagen synthesis consolidated. That the arena was fixed so rapidly, and the dissent of the
theory's own principal architects set aside, shows the configuration-space choice to have been a contingent turn
taken under the pressure of immediate predictive success and a positivist climate --- not a necessity forced by
the physics.

\section{RealQM as a three-dimensional primitive ontology}\label{sec:onto-realqm}

That a three-dimensional alternative was available, and not merely wished for, is the burden of this book: the
$N$ electrons carried by $N$ charge densities on ordinary space rather than by one amplitude on $3N$-space
(Chapter~\ref{ch:equation}). What the debate of \S\ref{sec:wfr-po} lets us add is its \emph{placement}: in those
terms RealQM is a \emph{primitive ontology in three dimensions}, its beables the charge densities in $\mathbb R^3$. Two placements keep it exact. It is not a hidden-variable theory --- it adds
no parameters beneath the amplitude (as de~Broglie--Bohm does) but \emph{replaces} the amplitude, as primary
object, with charge in three-space --- so the demand for ontology, often conflated with the demand for hidden
variables, is here met without any. And it stands to one side of the ontic/epistemic dispute altogether: since it
does not take the configuration-space $\Psi$ as its object, PBR's verdict on that object neither binds it nor is
contradicted by it --- what is real, for RealQM, is charge in three dimensions. RealQM is offered not as a
finished replacement --- its open problems, notably spin
exchange and the statistics of entangled measurements, are recorded in Chapter~\ref{ch:limits} --- but as an
existence proof that the road not taken in 1926 remains open.

\section{A practical corollary: quantum computing}\label{sec:onto-qc}

The arena problem is not purely interpretive; a technology can inherit it. The claimed advantage of quantum
computation rests on treating the joint state of $n$ qubits --- an amplitude in a $2^n$-dimensional space --- as
a physical resource on which operations act in parallel. This is the configuration-space object again, and
whether the exponential advantage is physically real depends on whether that object is a physical resource or a
formal device for the statistics of a system that is, in three-space, of modest size. On wave-function realism the
exponential resource is real and, decoherence permitting, exploitable; on a three-dimensional primitive ontology
there is no exponential thing in the device, and the scalable advantage lacks a substrate --- decoherence being
then not an incidental nuisance but the sign that the large superposition was never a robust physical thing. The
matter is not settled: small-scale devices demonstrably perform entangling operations and sampling tasks, and the
open question is the scalable, fault-tolerant regime, where a serious minority (Kalai, on grounds of noise) also
doubts that the barrier is merely technical. The philosophical observation is narrower and secure: the feasibility
of scalable quantum computing is entangled with what the exponentially large amplitude physically \emph{is} --- a
real resource or a formal device --- so the question of arena is not academic: it bears on where a great deal of
scientific effort and funding is now directed.

\section{Conclusion}\label{sec:onto-conclusion}

The debate over whether the wave function is ontic or epistemic can be left to the standard framework that
generated it. The point of this chapter is prior to that debate and independent of its outcome: whatever the
multi-dimensional $\Psi$ is taken to be, a theory whose fundamental object lives outside three-dimensional space
must recover the three-dimensional world by an added step it does not contain, and it is this that the century of
foundational trouble reflects. RealQM does not adjudicate the status of $\Psi$; it declines $\Psi$ as its object
and puts charge in three-space in its place, so that the three-dimensional facts are the theory's subject matter
rather than a reconstruction from it. The historical record shows the configuration-space choice to have been a
rapid and contested turn rather than a necessity, and the availability of a three-dimensional alternative shows
that the question closed in 1926 --- \emph{in what space is the physics of matter written?} --- was closed too
early, and too confidently.
